\def\stat{suchkov}



\def\tit{ЖИЗНЕННЫЙ ЦИКЛ КИБЕРАТАК\\ НА~УЯЗВИМОСТИ СИСТЕМ МАШИННОГО 

ОБУЧЕНИЯ}



\def\titkol{Жизненный цикл кибератак на~уязвимости систем машинного 

обучения}



\def\aut{А.\,П.~Сучков$^1$}



\def\autkol{А.\,П.~Сучков}



\titel{\tit}{\aut}{\autkol}{\titkol}



\index{Сучков А.\,П.}

\index{Suchkov A.\,P.}



%{\renewcommand{\thefootnote}{\fnsymbol{footnote}}

%\footnotetext[1]

%{Работа выполнялась с~использованием инфраструктуры Центра коллективного пользования 

%<<Высокопроизводительные вычисления и~большие данные>> (ЦКП <<Информатика>>) ФИЦ ИУ РАН 

%(г.~Москва).}} 



\renewcommand{\thefootnote}{\arabic{footnote}}

\footnotetext[1]{Федеральный исследовательский центр <<Информатика и~управление>> Российской 

академии наук, \mbox{ASuchkov@ipiran.ru}}





 \label{st\stat}



 \thispagestyle{headings}

 

\vspace*{-12pt}









\Abst{Рассматривается проблематика обеспечения необходимого уровня 

информационной безопас\-ности (ИБ) сис\-тем искусственного интеллекта (ИИ), что становится одним 

из ключевых факторов их широкого использования. На основе анализа данных 

о~совершенных кибератаках на известные сис\-те\-мы машинного обуче\-ния (СМО) 

проведено формирование полного набора стадий жизненного цик\-ла (ЖЦ) компьютерных атак 

на СМО, определены целевые уяз\-ви\-мости СМО для каждой стадии ЖЦ, 

сис\-те\-ма\-ти\-зи\-ро\-ва\-ны способы и~методы атак на уязвимости СМО, обозначены подходы 

к~обеспечению ИБ СМО.}



\KW{информационная безопас\-ность; сис\-те\-мы искусственного интеллекта; машинное 

обуче\-ние; виды защиты информации}



\DOI{10.14357\!/\!08696527240107}{QLGHEH}  



\vskip 10pt plus 9pt minus 6pt



 \thispagestyle{headings}

 

 \vspace*{-12pt}



\section{Введение}



\vspace*{-4pt}





    Системы, основанные на методах ИИ, 

переживают в~настоящий момент стадию бурного роста как по широте 

приложений, так и~по реальным достижениям по всем направлениям. Такие 

системы применяются в~области финансов, государственных услуг, 

в~медицине, транспорте и~многих других. Большой прогресс достигнут 

в~автоматизации распознавания образов, речи, в~развитии робототехники во 

многом благодаря существенному прогрессу СМО. Превратившись в~сис\-те\-мы высокой доступности и~массового 

использования системы ИИ все больше становятся объектом 

киберпреступлений, преследующих различные экономические, политические и~социальные цели и~опирающихся на высокотехнологичные методы, 

техники и~тактики.

    

     С целью обобщения опыта борьбы с~киберпреступлениями против 

СМО создана и~постоянно обновляется глобально доступная база знаний 

MITRE ATT\&CK$^\registered$~[1] о~тактике и~приемах противника, 

основанная на реальных наблюдениях. База знаний ATT\&CK используется в~качестве основы для разработки конкретных моделей угроз и~методов 

защиты информации в~государственном и~корпоративном секторе для 

создания эффективных продуктов и~услуг в~области ки\-бер\-безопас\-ности. 

База знаний ATT\&CK открыта и~доступна для бесплатного использования любым лицом 

или организацией.

     

     Российский аналог MITRE ATT\&CK~--- MITRE ATLAS$^{\mathrm{ТМ}}$ (Adversarial Threat Landscape for Artificial-Intelligence 

Systems)~[2], созданный факультетом безопас\-ности информационных 

технологий Университета ИТМО~[3],~--- это база знаний о~кейсах, тактиках 

и~техниках, применяемых злоумышленниками для атак на СМО. Она оформлена в~виде интерактивной матрицы, 

основанной на реальных наблюдениях, исследованиях комплексной 

имитации реальных атак с~целью оценки ки\-бер\-безопас\-ности сис\-тем 

и~групп по безопас\-ности. База знаний ATLAS создана по образцу фреймворка MITRE 

ATT\&CK$^\registered$, ее тактики и~техники дополняют матрицу ATT\&CK.

     

     База знаний ATLAS позволяет исследователям ориентироваться в~ландшафте угроз 

для СМО. Число уязвимостей в~них растет, и~их использование увеличивает 

фронт атак на существующие сис\-те\-мы. Матрица ATLAS была разработана 

для повышения осведомленности об угрозах ИБ 

СМО и~пред\-став\-ле\-ния их в~виде, удобном для исследователей безопас\-ности.

     %

     Базы знаний этих систем содержат более сотни примеров реальных и~имитированных кибератак на известные СМО, например: 

     \begin{itemize}

\item сбои служб машинного перевода~--- Google Translate, Bing Translator и~Systran Translate;

\item копирование модели GPT-2;

\item сбой службы Microsoft Azure;

\item обход системы ИИ в~классификаторе изображений Microsoft (on-edge);

\item выполнение вредоносного кода в~MathGPT с~помощью 

быстрого внедрения и~др.

\end{itemize}



     Техники, методы и~тактики атак в~базе знаний организованы в~виде 

матрицы по этапам их разработки и~выполнения. Интересно было бы 

упорядочить эти этапы по стадиям ЖЦ кибератак 

с~привязкой к~уязвимостям СМО и~методам их защиты. Итак, задачами 

исследования ставятся:

     \begin{enumerate}[(1)]

\item формирование полного набора стадий ЖЦ компьютерных атак на СМО;

\item определение целевых уязвимостей СМО для каждой стадии ЖЦ; 

\item систематизация способов и~методов атак на уязвимости СМО;

\item подходы к~обеспечению ИБ СМО.

\end{enumerate}



\section{Стадии жизненного цикла кибератак}



    Как всякая высокотехнологичная разработка, компьютерная атака на 

СМО проходит определенные стадии своего ЖЦ~--- от замысла, 

подготовки и~реализации до организованного отхода. Естественно, 

в~реальности не все эти стадии обязательны к~реализации: например, атака 

может носить спонтанный характер или иметь минимальную подготовку с~использованием готовых решений без тщательного мониторинга системы. 

Однако рассмотрим наиболее полный набор стадий ЖЦ для описания всех 

возможных тактик, способов и~техник, применяемых в~кибератаках. 

    

    Итак, на основе анализа данных~[1, 2] выделим 5 стадий ЖЦ 

кибератак с~указанием решаемых на этих стадиях задач.

    \begin{enumerate}[1.]

\item Целеполагание:

\begin{itemize}

\item[(а)] определение объекта атаки;

\item[(б)] формирование замысла атаки;

\item[(в)] планирование подготовки и~осуществления атаки.

\end{itemize}

\item Мониторинг объекта атаки:

    \begin{itemize}

\item[(а)] сбор доступной информации о СМО;

\item[(б)] исследование или сканирование системы жертвы;

\item[(в)] выявление уязвимостей объекта атаки;

\item[(г)] изучение существующего опыта атак уязвимостей известных 

моделей.

\end{itemize}

\item Подготовка ресурсов для атаки:

\begin{itemize}

\item[(а)] получение артефактов СМО;

\item[(б)] подготовка технической инфраструктуры для атаки;

\item[(в)] подготовка программных средств для атаки;

\item[(г)] публикация вредоносных данных для обучения;

\item[(д)] создание учетных записей к~различным сервисам.

\end{itemize}

\item Осуществление атаки: 

\begin{itemize}

\item[(а)] первоначальный доступ к~СМО;

\item[(б)] кража сведений о системе, имен учетных записей, паролей и~пр.;

\item[(в)] запуск вредоносного кода, встроенного в~артефакты или 

программное обеспечение (ПО);

\item[(г)] закрепление в~СМО;

\item[(д)] попытка получения доступа более высокого уровня;

\item[(е)] уклонение от защиты;

\item[(ж)] кража данных;

\item[(з)] манипулирование СМО и~данными, прерывание их работы или 

их уничтожение.

\end{itemize}

\item Предотвращение обнаружения и~идентификации субъекта атаки.

\end{enumerate}



\begin{table}\small  %1

\begin{center}

\tabcolsep=3pt

\begin{tabular}{|c|p{66mm}|p{50mm}|}

\multicolumn{3}{c}{Уязвимости СМО и~меры противодействия атакам на них}\\

\multicolumn{3}{c}{\ }\\[-6pt]

\hline

ЖЦ&\multicolumn{1}{c|}{Уязвимости}&\multicolumn{1}{c|}{Мероприятия по 

защите}\\

\hline

\tabcolsep=0pt\begin{tabular}{c}1а,\\ 1б,\\ 1в\end{tabular}&\multicolumn{1}{l|}{\tabcolsep=0pt\begin{tabular}{l}Популярность СМО.

Средний и~высокий\\ возможный ущерб\end{tabular}}&\multicolumn{1}{l|}{\tabcolsep=0pt\begin{tabular}{l}Учет степени повышенных рис-\\ ков при 

планировании меропри-\\ ятий по ИБ$^{*}$\end{tabular}}\\

\hline

2а&Доступность информации о научных исследованиях организации и~персоналий:\vspace*{-6pt}

\begin{itemize}

\addtolength{\itemsep}{-6pt}

\item репозитории препринтов;

\item технические блоги;

\item журналы и~материалы конференций;

\item информации с~сайта организации;

\item открытые репозитории аналогичных приложений (Google Play, iOS App Store, 

macOS App Store и~Microsoft Store)\vspace*{-6pt}

\end{itemize}&Анализ доступности критической информации о СМО и~учет при 

планировании системы защиты.\newline

Ограничение размещения в~открытом доступе критических данных$^{*}$\\

\hline

2б&Доступность данных о СМО:\vspace*{-6pt}

\begin{itemize}

\addtolength{\itemsep}{-6pt}

\item учетные данные;

\item адреса электронной почты;

\item имена сотрудников;

\item информация о сети;

\item информация об организации;

\item фишинг с~целью получения информации;

\item поиск в~закрытых источниках;

\item поиск в~открытых технических базах данных;

\item поиск открытых веб-сай\-тов\!/\!доменов\vspace*{-6pt}

\end{itemize}&Минимизация объема и~чувствительности данных, доступных внешним 

пользователям$^{*}$\\

\hline

2в&Уязвимости инфраструктуры:\vspace*{-6pt}

\begin{itemize}

\addtolength{\itemsep}{-6pt}

\item уязвимость серверов;

\item уязвимость целевых сред и~служб для эксплойта;

\item уязвимость программных средств;

\item уязвимость моделей\vspace*{-6pt}

\end{itemize}&Анализ и~учет уязвимостей инфраструктуры. Выявление аномалий в~потоках данных$^{*}$\\

\hline

2г&Наличие в~открытом доступе:\vspace*{-6pt}

\begin{itemize}

\addtolength{\itemsep}{-6pt}

\item исследования об уязвимостях распространенных моделей;

\item ранее полученные результаты, касающиеся нужного класса моделей, включая:\vspace*{-6pt}

\begin{itemize}

\addtolength{\itemsep}{-6pt}

\item публикации статей с~подробностями реализации успешной атаки;

\item ранее созданные реализации этих атак в~виде готовых программных средств\vspace*{-6pt}

\end{itemize}

\end{itemize}&Поиск в~открытом доступе данных об успешных атаках на выбранную 

модель.\newline

Учет степени повышенных рисков при планировании мероприятий по ИБ$^{*}$\\

\hline

\multicolumn{3}{r}{\textit{Продолжение таблицы на с.}~\pageref{s1}}

\end{tabular}

\end{center}

\end{table}





    Упорядочим данные баз знаний~[1, 2] в~виде таблицы. 

Следует отметить, что в~этих базах приведено около двухсот техник и~тактик 

атак, поэтому в~таб\-ли\-це в~ряде случаев указано их зарегистрированное 

число и~перечислены наиболее представительные типы.





\begin{table}\small  %2

\label{s1}

\begin{center}

\tabcolsep=3pt

\begin{tabular}{|c|p{66mm}|p{50mm}|}

\multicolumn{3}{c}{Уязвимости СМО и~меры противодействия атакам на них (\textit{продолжение})}\\

\multicolumn{3}{c}{\ }\\[-6pt]

\hline

ЖЦ&\multicolumn{1}{c|}{Уязвимости}&\multicolumn{1}{c|}{Мероприятия по 

защите}\\

\hline

3а&Наличие в~общедоступных источниках, включая облачные хранилища, 

общедоступных сервисов, а~также ПО репозиториев данных, 

артефактов машинного обучения:\vspace*{-6pt}

\begin{itemize}

\addtolength{\itemsep}{-6pt}

\item программный стек, используемый для обучения и~развертывания моделей; 

\item данные обучения и~тестирования;

\item конфигурации и~параметры мо\-делей\vspace*{-6pt}

\end{itemize}&Поиск в~общедоступных источниках данных об артефактах СМО. \newline

Учет степени повышенных рисков при планировании мероприятий по ИБ\\

\hline

3б&Уязвимости доступной инфраструктуры с~использованием:\vspace*{-6pt} 

\begin{itemize}

\addtolength{\itemsep}{-6pt}

\item доменных имен (похожие, устаревшие и~др.);

\item DNS-серверов;

\item виртуальных частных серверов;

\item своих серверов;

\item ботнетов;

\item популярных веб-сер\-ви\-сов;

\item бессерверного подхода;

\item вредоносной рекламы\vspace*{-6pt}

\end{itemize}&Регистрация доменов с~похожими именами.\newline

Сканирование интернета на предмет выявления серверов, приобретенных 

злоумышленниками.\newline

Возможен поиск уникальных характеристик, связанных с~ПО

противника, если они известны.\newline

Средства блокировки рекламы для предот\-вра\-ще\-ния выполнения вредоносного кода$^{*}$\\

\hline

3в&Возможности получения ПО:\vspace*{-6pt}

\begin{itemize}

\addtolength{\itemsep}{-6pt}

\item бесплатное скачивание ПО с~открытым кодом;

\item покупка ПО; 

\item кража ПО и/или лицензии на ПО у~сторонних организаций;

\item взлом пробных версий;

\item разработка собственных версий ПО\vspace*{-6pt}

\end{itemize}&Радикально~--- отказ от использования общедоступного недоверенного 

ПО, иначе см.$^{*}$\\

\hline

3г&Внедренные уязвимости в~модели машинного обучения:\vspace*{-6pt}

\begin{itemize}

\addtolength{\itemsep}{-6pt}

\item подготовка и~опубликование вредоносных данных для обучения;

\item отравление существующих данных для обучения;

\item создание возможности активации уязвимости образцами данных  

с~бэк\-дор-триг\-ге\-ром\vspace*{-6pt} 

\end{itemize}&Радикально ~--- отказ от использования общедоступных недоверенных 

обуча\-ющих данных.\newline

Мероприятия по нормализации обуча\-ющих наборов данных и~по повышению 

устойчивости модели СМО\\

\hline

\multicolumn{3}{r}{\textit{Продолжение таблицы на с.}~\pageref{s2}}

\end{tabular}

\end{center}

\end{table}



\begin{table}\small  %3

\label{s2}

\begin{center}

\tabcolsep=3pt

\begin{tabular}{|c|p{64mm}|p{52mm}|}

\multicolumn{3}{c}{Уязвимости СМО и~меры противодействия атакам на них (\textit{продолжение})}\\

\multicolumn{3}{c}{\ }\\[-6pt]

\hline

ЖЦ&\multicolumn{1}{c|}{Уязвимости}&\multicolumn{1}{c|}{Мероприятия по 

защите}\\

\hline

3д&Уязвимости, связанные с~доступом к~сервисам для использования в~целевых атаках, 

чтобы получить доступ к~необходимым ресурсам для постановки атак или для выдачи 

себя за жертву:\vspace*{-6pt}

\begin{itemize}

\addtolength{\itemsep}{-6pt}

\item создание и~поддержание учетных записей в~сервисах, которые могут быть 

использованы во время целевых атак~--- социальные сети, электронная почта, сайты 

и~т.\,п.;

\item создание и~использование он\-лайн-пер\-со\-нажа\vspace*{-6pt}

\end{itemize}&Анализ содержания сетевого трафика.\newline

Мониторинг активности в~социальных сетях, связанных с~организацией СМО$^{*}$\\

\hline

4а&Уязвимости в~системе доступа по различным векторам входа, используемые для 

закрепления в~системе\vspace*{-6pt}

\begin{itemize}

\addtolength{\itemsep}{-6pt}

\item компрометация цепочек поставок аппаратно-программных средств, ПО, данных, 

моделей;

\item доступ к~модели СМО через API, доступ к~данным из физической среды, полный 

доступ к~модели, косвенно через взаимодействие с~продуктом или сервисом, 

использующим машинное обуче\-ние;

\item использование существующих учетных записей;

\item обход систем обнаружения вредоносного ПО или сетевого сканирования, 

основанных на машинном обучении;

\item репликация через съемные носители;

\item фишинг\vspace*{-6pt}

\end{itemize}&Внедрение процесса управ\-ле\-ния ис\-прав\-ле\-ни\-ями для проверки 

не\-ис\-поль\-зу\-емых зависимостей.\newline

Внедрение непрерывного мониторинга источников уязвимостей.\newline

Использование многофакторной идентификации.\newline

Сегментация сети.\newline

Надлежащее управление учетными записями и~разрешениями, используемыми сторонами в~доверительных отношениях.\newline

Ограничение использования 

 USB-устройств и~съемных носителей в~сети\\

\hline

4б&Уязвимости в~политике ведения учетных записей:\vspace*{-6pt}

\begin{itemize}

\addtolength{\itemsep}{-6pt}

\item устаревший порядок использования учетных записей;

\item недостатки в~политике ведения паролей\vspace*{-6pt}

\end{itemize}

&Использование политики условного доступа для блокировки входов с~устройств, не 

соответствующих требованиям, или с~IP-адресов за пределами определенных 

организацией диапазонов.\newline

Ключи должны периодически об\-нов\-лять\-ся и~надлежащим образом защищаться.\newline

Отслеживание необычного поведения учетных записей\\

\hline

\multicolumn{3}{r}{\textit{Продолжение таблицы на с.}~\pageref{s3}}

\end{tabular}

\end{center}

\end{table}







\begin{table}\small  %4

\label{s3}

\begin{center}

\tabcolsep=3pt

\begin{tabular}{|c|p{58mm}|p{58mm}|}

\multicolumn{3}{c}{Уязвимости СМО и~меры противодействия атакам на них (\textit{продолжение})}\\

\multicolumn{3}{c}{\ }\\[-6pt]

\hline

ЖЦ&\multicolumn{1}{c|}{Уязвимости}&\multicolumn{1}{c|}{Мероприятия по 

защите}\\

\hline

4в&Использование уязвимостей с~возможностями запуска вредоносного кода:\vspace*{-6pt}

\begin{itemize}

\addtolength{\itemsep}{-6pt}

\item запуск пользователем с~умыслом или по невнимательности;

\item запуск с~помощью артефактов машинного обуче\-ния;

\item использование инструментов удаленного доступа для запуска сценария\vspace*{-6pt}

\end{itemize}&Существует более~40~видов возможностей запуска вредоносного кода и~техник защиты от них~[1], в~частности администрирование:\vspace*{-6pt}

\begin{itemize}

\addtolength{\itemsep}{-6pt}

\item облачных сервисов;

\item интерпретаторов команд и~сценариев;

\item служб контейнеров;

\item клиентских приложений;

\item механизмов межпроцессной связи;

\item создания собственных API;

\item планирования задач; 

\item функций пользователей\vspace*{-6pt}

\end{itemize}\\

\hline

4г&Уязвимости системы доступа при перезапусках, изменении учетных данных и~других 

прерываниях, которые могут привести к~прекращению доступа:\vspace*{-6pt}

\begin{itemize}

\addtolength{\itemsep}{-6pt}

\item манипулирование сроками действия паролей;

\item использование фоновых задач;

\item системы автозапуска;

\item сценариев инициализации загрузки и~входа в~систему;

\item внешних удаленных сервисов;

\item процессов аутоидентификации\vspace*{-6pt} 

\end{itemize}&Существует более 20 видов возможностей закрепления и~техник 

защиты от них~[1], в~частности администрирование:\vspace*{-6pt}

\begin{itemize}

\addtolength{\itemsep}{-6pt}

\item сроков действия паролей;

\item фоновых задач;

\item системы автозапуска;

\item сценариев инициализации загрузки и~входа в~систему;

\item внешних удаленных сервисов;

\item процессов аутоидентификации\vspace*{-6pt}

\end{itemize}\\

\hline

4д&Уязвимости в~разграничении прав доступа на уровне:\vspace*{-6pt}

\begin{itemize}

\addtolength{\itemsep}{-6pt}

\item системном~--- корневой уровень;

\item локального администратора;

\item учетной записи пользователя с~правами администратора;

\item учетных записей пользователей, имеющих доступ к~определенной сис\-те\-ме или 

выполняющих определенную функцию\vspace*{-6pt}

\end{itemize}&Существует более 10~видов возможностей атак и~техник защиты от 

них~[1], в~частности администрирование:\vspace*{-6pt}

\begin{itemize}

\addtolength{\itemsep}{-6pt}

\item аудита операционной сис\-темы;    

\item выполнения приложений;    

\item конфигурации операционной сис\-темы;    

\item управления привилегированными учетными записями;  

\item прав доступа к~файлам и~каталогам\vspace*{-6pt}

\end{itemize}\\

\hline

4е&Методы, используемые для уклонения от защиты, включают:\vspace*{-6pt} 

\begin{itemize}

\addtolength{\itemsep}{-6pt}

\item удаление/отклю\-че\-ние ПО без\-опас\-ности;

\item запутывание/шифрование данных и~сценариев;

\item использование доверенных процессов и~злоупотребление ими\vspace*{-6pt}

\end{itemize}&Существует более 40~видов возможностей атак и~техник защиты от них~[1], в~частности администрирование:\vspace*{-6pt} 

\begin{itemize}

\addtolength{\itemsep}{-6pt}

\item системы ПО ИБ;

\item системы архивирования и~восстановления;

\item доверенных процессов\vspace*{-6pt}

\end{itemize}\\

\hline

\multicolumn{3}{r}{\textit{Продолжение таблицы на с.}~\pageref{s4}}

\end{tabular}

\end{center}

\vspace*{-1.83685pt}

\end{table}



\begin{table}\small 

\label{s4}

\begin{center}

\tabcolsep=3pt

\begin{tabular}{|c|p{58mm}|p{58mm}|}

\multicolumn{3}{c}{Уязвимости СМО и~меры противодействия атакам на них (\textit{продолжение})}\\

\multicolumn{3}{c}{\ }\\[-6pt]

\hline

ЖЦ&\multicolumn{1}{c|}{Уязвимости}&\multicolumn{1}{c|}{Мероприятия по 

защите}\\

\hline

4ж&Уязвимости в~системе доступа к~модели машинного обуче\-ния, дающие возможности:\vspace*{-6pt}

\begin{itemize}

\addtolength{\itemsep}{-6pt}

\item получения данных по своему каналу управ\-ле\-ния или альтернативному каналу;

\item установления ограничений на размер передачи; 

\item кражи артефактов машинного  

обуче\-ния;

\item вывода о принадлежности образца данных к~обуча\-юще\-му набору, что считается 

нарушением конфиденциальности;

\item извлечение копии модели\vspace*{-6pt}

\end{itemize}&Этот тип атак не может быть легко устранен с~помощью превентивных мер 

контроля, поскольку он основан на злоупотреблении системными функциями.\newline

Использование системы обнаружения и~предотвращения сетевых вторжений. \newline

Настройка сетевого брандмауэра, чтобы разрешать вход в~сеть и~выход из нее только 

необходимым портам.\newline

Строгий контроль идентификации и~управления доступом\\

\hline

4з&Уязвимости СМО и~данных, позволяющие прерывать работу или их уничтожить 

путем:\vspace*{-6pt} 

\begin{itemize}

\addtolength{\itemsep}{-6pt}

\item создания состязательных данных, которые нарушают корректную идентификацию 

содержимого данных моделью;

\item создания потока запросов с~целью ухудшить или остановить работу всего сервиса;

\item ухудшения про\-из\-во\-ди\-тель\-ности целевой модели, направляя бесполезные запросы 

или вычислительно сложные входные данные на различные сервисы машинного 

обучения, чтобы увеличить затраты на работу этих сервисов;

\item извлечения артефактов машинного обучения, чтобы украсть интеллектуальную 

собственность и~причинить экономический вред;

\item засорение мусорными данными, что приводит к~увеличению числа ложных 

обнаружений\vspace*{-6pt}

\end{itemize}&Существует 14~видов атак и~техник защиты от них~[1], в~частности:\vspace*{-6pt}

\begin{itemize}

\addtolength{\itemsep}{-6pt}

\item внедрение планов аварийного восстановления информационных технологий, содержащих процедуры 

регулярного создания резервных копий данных, которые могут использоваться для 

восстановления данных организации;

\item фильтрация пограничного трафика, блокируя адреса источников, на которые 

направлена атака, блокируя порты, на которые нацелены, или блокируя протоколы, 

используемые для транспортировки;

\item обучение пользователей распознавать методы социальной инженерии, используемые 

для совершения финансовых краж;

\item отслеживание выполняемых команд и~аргументов двоичных файлов, участвующих в~выключении или перезагрузке систем\vspace*{-6pt}

\end{itemize}\\

\hline

\multicolumn{3}{r}{\textit{Окончание таблицы на с.}~\pageref{s5}}

\end{tabular}

\end{center}

\end{table}





%\vspace*{-9pt}





\section{Подходы к~защите информации для систем машинного обучения}

     

     К основным видам защиты информации для СМО можно отнести 

(ГОСТ Р~50922-2006):

     \begin{itemize}

\item техническую защиту информации;

\item криптографическую защиту информации, включая обуча\-ющие 

данные;

\item физическую защиту информации; 

\item правовую защиту информации;

\end{itemize}



\begin{table}\small  %5

\label{s5}

\begin{center}

\tabcolsep=3pt

\begin{tabular}{|c|p{70mm}|p{46mm}|}

\multicolumn{3}{c}{Уязвимости СМО и~меры противодействия атакам на них (\textit{окончание})}\\

\multicolumn{3}{c}{\ }\\[-6pt]

\hline

ЖЦ&\multicolumn{1}{c|}{Уязвимости}&\multicolumn{1}{c|}{Мероприятия по 

защите}\\

\hline

5&Уязвимости в~системах безопас\-ности, в~том чис\-ле, основанных на методах машинного обуче\-ния для 

предотвращения отслеживания~[4]:\vspace*{-6pt}

\begin{itemize}

\addtolength{\itemsep}{-6pt}

\item использование обычных сервисов, таких как Google или Twitter, что облегчает 

злоумышленникам скрытие в~ожидаемом шуме; 

\item маскировка путем переименования инструментов и~двоичных файлов в~соответствии с~файлами и~программами на взломанном устройстве;

\item отключение регистрации событий с~помощью AUDITPOL перед практическими 

действиями с~клавиатуры и~включение обратно после них;

\item создание правил брандмауэра для минимизации исходящих пакетов для 

определенных протоколов перед запуском действий по пе\-ре\-чис\-ле\-нию сетей, которые 

потом удалялись;

\item отключение служб безопас\-ности на целевых хостах, использование временных 

меток для изменения временных меток артефактов, а также процедур и~инструментов 

очистки, чтобы препятствовать обнаружению вредоносных внедрений DLL в~уязвимых 

средах\vspace*{-6pt}

\end{itemize}&Необходимо реализовывать комплексные меры обеспечения ИБ, к~одной из 

первоочередных можно отнести процессы мониторинга и~реагирования на 

киберинциденты. В~случае недостатка собственных сил и~средств на рынке РФ 

существует достаточное число провайдеров сервисов безопас\-ности~[4]\\

\hline

\multicolumn{3}{p{128mm}}{\footnotesize \hspace*{3mm}$^{*}$Кибератаки на эти уязвимости не 

могут быть легко предотвращены с~помощью превентивных мер защиты, 

поскольку они основаны на подходах, выходящих за рамки средств защиты и~контроля предприятия.}

\end{tabular}

\end{center}

\vspace*{-12pt}

\end{table}



К основным видам защиты СМО необходимо добавить~[5, 6]:

\begin{itemize}

\item обеспечение системной и~функциональной безопас\-ности;

\item повышение устойчивости СМО.

\end{itemize}

     

     Техническая и~физическая защита, обеспечивающая 

конфиденциальность, целостность и~доступность информации, 

осуществляется на основе существующих  

нор\-ма\-тив\-но-ме\-то\-ди\-че\-ских документов и~не отличается от 

комплекса мероприятий по обеспечению безопас\-ности информации для 

любых автоматизированных информационных сис\-тем. При этом следует 

обратить внимание на то, что на начальных этапах ЖЦ кибератак эти виды 

защиты затруднены, так как выходят за пределы компетенции предприятия 

и~могут быть реализованы только путем ограничения взаимодействия СМО 

с~внешними системами и~сервисами и~выполнения персоналом 

определенных обязательств.

  %   

     Другая особенность обеспечения ИБ СМО 

кроется в~их специфических свойствах. Одна из основных проб\-лем 

обеспечения ИБ СМО на основе технологии 

машинного обучения заключается в~том, что создание и~функционирование 

таких систем невозможно без разнообразных процедур обучения, 

до\-обуче\-ния и~са\-мо\-обуче\-ния с~использованием зачастую огромных 

массивов данных из различных источников~--- баз данных, баз знаний, 

экспертных данных, коллекций изображений и~аудиозаписей, сред 

сенсорного и~имитационного моделирования, <<цифровых двойников>>. 

Очень часто источники этих данных не входят в~число доверенных и~формируются в~открытом доступе с~целью тестирования новых подходов 

к~созданию СМО самыми разными авторами и~научным сообществом в~целом. 

     

     Еще одна специфическая особенность~--- фактор применения 

интеллектуальных автономных и~мобильных СМО, что предъявляет 

специфические требования к~сохранению конфиденциальности информации и~информационных процессов в~смысле расширения спектра угроз ИБ 

(возможность извлечения и~инверсии моделей и~т.\,п.).

     

     С точки зрения информационных процессов, влияющих на 

ИБ СМО, к~одним из основных специфических 

угроз следует отнести процедуры обучения систем ИИ, в~ходе которых 

настраивается и~может меняться не только информационная база, но 

и~алгоритмы интеллектуальных выводов. Эта проб\-ле\-ма устой\-чи\-вости 

СМО состоит в~том, что даже небольшие изменения в~обучающих наборах 

данных могут привести к~кардинальным изменениям интеллектуальных 

способностей таких сис\-тем, вплоть до их существенного нарушения. 

Важно также очень трудно проверяемое качество обуча\-ющих данных 

(репрезентативность, статистическая несмещенность, достаточный объем).

     

     Таким образом, с~точки зрения ИБ СМО 

данные свойства таких систем предоставляют злоумышленникам более 

широкий спектр уязвимостей для вредоносного воздействия. Специфика 

защиты информации и~СМО может потребовать 

дополнительных мер на стадии их проектирования и~разработки, таких как 

изменение гиперпараметров моделей, применения процедур верификации 

обучающих данных, создание интеллектуальных подсистем тестирования на 

предмет устойчивости к~вредоносным воздействиям. Это осуществляется 

в~рамках обеспечения \textit{сис\-тем\-ной} безопас\-ности путем 

определения целевых областей применения СМО и~ее критических 

элементов, формирования реализуемых системных требований на основе 

анализа эксплуатационных и~технологических рисков.



     

\section{Заключение}



\vspace*{-4pt}



\noindent

\begin{enumerate}[1.]

\item Системы машинного обучения относятся к~одному из видов автоматизированных 

информационных систем, и~к ним должны применяться все су\-щест\-ву\-ющие 

виды защиты информации (правовая, техническая, криптографическая, 

физическая, функциональная, системная), регламентированные 

существующей нор\-ма\-тив\-но-ме\-то\-ди\-че\-ской базой.

\item Основные специфические уязвимости СМО несут процедуры обуче\-ния 

сис\-тем ИИ, в~ходе которых настраивается и~может меняться не только 

информационная база, но и~алгоритмы интеллектуальных выводов.

\item Предложена классификация уязвимостей СМО и~мер защиты 

информации по стадиям ЖЦ кибератак. 

\item Обсуждены специфические вопросы защиты информации СМО и~ряд 

мероприятий по обеспечению ИБ этих сис\-тем. 

\end{enumerate}



\vspace*{-8pt}



{\small\frenchspacing

{%\baselineskip=10.8pt

\addcontentsline{toc}{section}{Литература}

\begin{thebibliography}{9}



\bibitem{1-suc}

MITRE ATT\&CK$^\registered$. {\sf https://attack.mitre.org.}

\bibitem{2-suc}

MITRE ATLAS$^{\mathrm{TM}}$ (Adversarial Threat Landscape for  

Artificial-Intelligence Systems). {\sf https://atlas.securityhub.ru}.

\bibitem{3-suc}

Университет ИТМО. {\sf https://itmo.ru/?ysclid=lqbyvy76bx557068225}.

\bibitem{4-suc}

Microsoft рассказала, как хакеры избежали обнаружения при атаке SolarWinds. {\sf 

https://habr.com/ru/search/?q=уклонение\%20от\%20обнаружения\&target\_type=\linebreak posts\&order=relevance}. 



\bibitem{6-suc}

\Au{Гаврилов В.\,Е., Зацаринный А.\,А.} Исследование проблем  

нор\-ма\-тив\-но-ме\-то\-ди\-че\-ско\-го регулирования в~об\-ласти информационной 

безопас\-ности процессов создания и~внед\-ре\-ния информационных технологий, 

разрабатываемых в~рамках программы <<Цифровая экономика>>~/\!\!/ Вестник 

Воронежского института ФСИН России, 2020. №\,3. C.~30--37. EDN: HNRNCW.

\bibitem{5-suc}

\Au{Зацаринный А.\,А., Сучков А.\,П.} Некоторые подходы к~анализу факторов, влияющих 

на информационную безопас\-ность сис\-тем искусственного интеллекта~/\!\!/ Сис\-те\-мы 

и~средства информатики, 2023. Т.~33. №\,3. С.~95--107. doi: 

10.14357\!/ 08696527230308. EDN: QNXZBN.

   \end{thebibliography}

} }







\vspace*{-9pt}



\hfill{\small\textit{Поступила в~редакцию 20.12.23}}



\vspace*{5pt}



\hrule



\vspace*{2pt}





\hrule



\vspace*{6pt}



%\newpage



%\vspace*{-28pt}



\def\tit{THE LIFECYCLE OF~CYBERATTACKS\\ ON~MACHINE LEARNING 

SYSTEM VULNERABILITIES\\[-5pt]}







\def\titkol{The lifecycle of cyberattacks on~machine learning 

system vulnerabilities}



\def\aut{A.\,P.~Suchkov\\[-5pt]}



\def\autkol{A.\,P.~Suchkov}



\titel{\tit}{\aut}{\autkol}{\titkol}



\vspace*{-15pt}



\noindent

Federal Research Center ``Computer Science and Control'' of the Russian Academy

of Sciences,   44-2~Vavilov Str., Moscow 119133, Russian Federation





\def\leftfootline{\small{\textbf{\thepage}

\hfill Sistemy i Sredstva Informatiki~---

Systems and Means of Informatics 2024 vol~34 no\,1}

}%

 \def\rightfootline{\small{Sistemy i Sredstva Informatiki~---

 Systems and Means of Informatics 2024 vol~34 no\,1

\hfill \textbf{\thepage}}}



%\vspace*{1pt}





\Abste{The article discusses the problems of ensuring the necessary level of 

information security of artificial intelligence systems which is becoming one of 

the\linebreak\vspace*{-12pt}}



\Abstend{key factors in their widespread use. Based on the analysis of data on committed 

cyberattacks on well-known machine learning systems (MLS), a complete set of 

stages of the life cycle of computer attacks on MLS was formed, target MLS 

vulnerabilities for each stage of the life cycle were determined, methods of attacks 

on MLS vulnerabilities were systematized, and approaches to ensuring information 

security of the system were outlined.}



%\vspace*{-10pt}



\KWE{information security; artificial intelligence systems; machine learning; 

types of information protection}



%\vspace*{-10pt}



\DOI{10.14357\!/\!08696527240107}{QLGHEH}  



%\vspace*{-14pt}





%\Ack

%\vspace*{-3pt}

%\noindent





%\vspace*{-3pt}



\renewcommand{\bibname}{\protect\rmfamily References}

%\renewcommand{\bibname}{\large\protect\rm References}



{\small\frenchspacing

{ %\baselineskip=10.8pt

\addcontentsline{toc}{section}{References}

\begin{thebibliography}{9}



%\vspace*{-7pt}



\bibitem{1-suc-1}

MITRE ATT\&CK$^\registered$. Available at: {\sf https://attack.mitre.org} 

(accessed February~20, 2024).

\bibitem{2-suc-1}

MITRE ATLAS$^{{\mathrm{TM}}}$ (Adversarial Threat Landscape for 

Artificial-Intelligence Systems). Available at: {\sf https://atlas.securityhub.ru} (accessed 

February 20, 2024).

 \bibitem{3-suc-1}

Universitet ITMO [ITMO University]. Available at: {\sf 

https://itmo.ru/?ysclid=\linebreak lqbyvy76bx557068225} (accessed February~20, 2024).

\bibitem{4-suc-1}

Microsoft rasskazala, kak khakery izbezhali obnaruzheniya pri atake SolarWinds 

[Microsoft told how hackers evaded detection in the SolarWinds attack]. Available 

at: {\sf https://habr.com/ru/news/538476/} (accessed February~20, 2024).



\bibitem{5-suc-1} %6

\Aue{Zatsarinnyy, A.\,A., and A.\,P.~Suchkov.} 2023. Nekotorye podkhody 

k~analizu faktorov, vliyayushchikh na informatsionnuyu bezopas\-nost' sis\-tem 

iskusstvennogo intellekta [Some approaches to the analysis of factors affecting the 

information security of artificial intelligence systems]. \textit{Sistemy i~Sredstva 

Informatiki~--- Systems and Means of Informatics} 33(3):95--107. doi: 

10.14357/08696527230308. EDN: QNXZBN.

\bibitem{6-suc-1} %5

\Aue{Gavrilov, V.\,E., and A.\,A.~Zatsarinnyy.} 2020. Issledovanie problem 

normativno-metodicheskogo regulirovaniya v~oblasti informatsionnoy 

bezopasnosti protsessov so\-zda\-niya i~vnedreniya informatsionnykh tekhnologiy, 

razrabatyvaemykh v~ramkakh pro\-gram\-my ``Tsifrovaya ekonomika'' [Study of the 

problems of normative-methodical management in the field of information security 

processes for the creation and implementation of information technologies 

developed in the framework of the ``Digital economy'']. \textit{Vestnik 

Voronezhskogo instituta FSIN Rossii} [Bulletin of the Voronezh Institute of the 

Federal Penitentiary Service of Russia] 3:30--37. EDN: HNRNCW.

\end{thebibliography}

} }



\vspace*{-6pt}



\hfill{\small\textit{Received December 20, 2023}} 



\vspace*{-14pt} 



\Contrl



\vspace*{-3pt}



\noindent

\textbf{Suchkov Alexander P.} (b.\ 1954)~--- Doctor of Science in technology, 

leading scientist, Federal Research Center ``Computer Science and Control'' of the 

Russian Academy of Sciences, 44-2~Vavilov Str., Moscow 119333, Russian 

Federation; \mbox{ASuchkov@frccsc.ru}







  

\label{end\stat}



\renewcommand{\bibname}{\protect\rm Литература} 

